\documentclass{article}
\input{~/studies/preamble}

\title{Gaussian Mixture Model}
\author{Hyoung Joon, Kim}
\date{\today}

\begin{document}
\maketitle
\tableofcontents
\newpage

\section{Introduction}
    This document explains \textbf{Gaussian mixture model} (GMM), which is literally a mixture of Gaussian distribution.
    GMM is mainly used for clustering like k-means clustering, but the main difference is that it provides us the probability of being member of each clusters.

\section{Main Idea}
    \subsection{Introduction}
        Gaussian mixture model can be simply illustrated as a linear superposition of Gaussian distributions;
        \[
        p(\textbf{x})=\sum_{k=1}^{K}\pi_k\mathcal{N}(\textbf{x}|\bm\mu_k, \bm\Sigma_k),\ 
        \text{where $0 \leq \pi_k \leq 1$ and $\sum_{k=1}^{K}\pi_k = 1$}.
        \]
        However, the GMM can be expressed in terms of discrete \emph{latent variables}, and its more useful when interpreted in this way.
    \subsection{Latent Variable}
        \begin{defn}{Gaussian Mixture Model}{gaussMixMod}
            Let \textbf{z} be a $K$-dimensional binary random variable having 1-of-$K$ representation i.e.,
            $z_k \in \{0,1\}$ and $\sum_k z_k=1$. If the random variable \textbf{x} have conditional probability distribution
            \[
            p(\mathbf{x}|z_k=1) = \mathcal{N}(\mathbf{x}|\bm\mu_k,\bm\Sigma_k)
            \]
            with
            \[
            p(z_k=1)=\pi_k,\text{ where } 0 \leq \pi_k \leq 1 \text{ and } \sum_{k=1}^{K}\pi_k=1
            \]
            \textbf{x} follows mixture of Gaussians.
        \end{defn}

        We can derive the linear superposition representation of Definition \ref{defn:gaussMixMod} when we calculate the marginal distribution $p(\textbf{x})$.
        The advantage of this representation is that we can directly work with joint distribution $p(\textbf{x}, \textbf{z})$, and we can evaluate the \emph{responsibility} $\gamma(z_{nk})$ of 
        each latent variable $z_k$ about data point $\textbf{x}_n$, which is defined as
        \begin{align*}
        \gamma(z_{nk})&=p(z_k=1|\textbf{x}_n) \\
        &=\frac{p(z_k=1)p(\textbf{x}_n|z_k=1)}{\sum_{j=1}^{K}p(z_j=1)p(\textbf{x}_n|z_j=1)} \\
        &=\frac{\pi_k\mathcal{N}(\textbf{x}_n|\bm\mu_k,\bm\Sigma_k)}{\sum_{j=1}^{K}\mathcal{N}(\textbf{x}_n|\bm\mu_j,\bm\Sigma_j)}
        \end{align*}
    
    \subsection{Likelihood Function}
        The likelihood function of a mixture of Gaussians is given as following:
        \begin{thm}{Likelihood Function of GMM}{likeliGMM}
            Let \textbf{X} be a given data set in which the $n$-th row is given by $\textbf{x}_n^{\top}$. Using the notations of Definition \ref{defn:gaussMixMod}, the log-likelihood is given as
            \[
            \log{p(\textbf{X}|\bm\pi,\bm\mu,\bm\Sigma)}=\sum_{n=1}^{N}\log{\left\{ \sum_{k=1}^{K}\pi_k\mathcal{N}(\textbf{x}_n|\bm\mu_k,\bm\Sigma_k) \right\}}
            \]
        \end{thm}
        The maximum likelihood, unlike simple Gaussian model, cannot be expressed in closed form. To maximize the log-likelihood function, we can use \textbf{EM} algorithm, which is usedful for
        optimization with latent variables. General EM algorithm explanation is given in this \href{run:../../inference/em_algorithm/em_algorithm.pdf}{document}.

        Before discussing about the EM algorithm for GMM, we show that there is a problem with the maximum likelihood framework when it is applied to GMM.
        Suppose that the $j$-th component $\bm\mu_j$ of $\bm\mu$ exactly equal to one of the data points, i.e., $\bm\mu_j=\textbf{x}_n$ for some $n$. Then, the
        likelihood of the GMM will diverge to infinity, since assuming that the variance is isotropic,
        $\mathcal{N}(\textbf{x}_n|\bm\mu_j, \sigma_j^2\textbf{I})=\frac{1}{\sqrt{2\pi}\sigma_j} \rightarrow \infty$ as $\sigma_j \rightarrow 0$, when maximizing
        likelihood. This can be resolved heuristically by detecting when a Gaussian component is collapsing, and then randomly assigning values to mean and resetting variance to some large value.
        Or we can set the prior on parameters and try to get MAP estimate instead of MLE.

    \subsection{EM algorithm for GMM}
        \begin{algorithm}[H]
            \DontPrintSemicolon
            \SetKwInOut{Input}{Input}
            \SetKwInOut{Output}{Output}
            \SetArgSty{textnormal}

            \caption{EM algorithm for a Gaussian mixture model}

            \Input{initial model parameters $\{\bm\mu_k\},\{\bm\Sigma_k\},\{\pi_k\}$ \newline
            data set $\{\textbf{x}_1,\dots,\textbf{x}_n\}$}
            \Output{final model parameters $\{\bm\mu_k\},\{\bm\Sigma_k\},\{\pi_k\}$}

            \BlankLine
            \hrule

            \Repeat{convergence}
            {
                \tcp{E step}
                \For{$n \in \{1,\dots,N\}$}
                {
                    \For{$k \in \{1,\dots,K\}$}
                    {
                        $\gamma(z_{nk}) \leftarrow \dfrac{\pi_k\mathcal{N}(\textbf{x}_n|\bm\mu_k,\bm\Sigma_k)}{\sum_{j=1}^{K}\pi_j\mathcal{N}(\textbf{x}_n|\bm\mu_j,\bm\Sigma_j)}$\;
                    }
                }
                \tcp{M step}
                \For{$k \in \{1,\dots,K\}$}
                {
                    $N_k \leftarrow \sum_{n=1}^{N}\gamma(z_{nk})$\;
                    $\bm\mu_k \leftarrow \frac{1}{N_k}\sum_{n=1}^{N}\gamma(z_{nk})\textbf{x}_n$\;
                    $\bm\Sigma_k \leftarrow \frac{1}{N_k}\sum_{n=1}^{N}\gamma(z_{nk})(\textbf{x}_n-\bm\mu_k)(\textbf x_n-\bm\mu_k)^{\top}$\;
                    $\pi_k \leftarrow \frac{N_k}{N}$\;
                }
                \tcp{log likelihood}
                $\mathcal{L} \leftarrow \sum_{n=1}^{N}\log{\left\{ \sum_{k=1}^{K}\pi_k\mathcal{N}(\textbf{x}_n|\bm\mu_k,\bm\Sigma_k) \right\}}$\;
            }
            \Return{$\{\bm\mu_k\},\{\bm\Sigma_k\},\{\pi_k\}$}

        \end{algorithm}
\end{document}